\section*{Summary}\label{summary}
\addcontentsline{toc}{section}{Summary}

Alexandre Bernardino holds a PhD in Electrical and Computer Engineering
from \emph{Instituto Superior Técnico} (IST), \emph{University of
Lisbon} (UL) (2004). He is currently an Associate Professor with tenure
at the scientific area of \emph{Systems, Decision and Control} (ACSDC)
of the \emph{Electrical and Computer Engineering Department} (DEEC). He
is a researcher of the \emph{Institute for Systems and Robotics of IST}
(ISR/IST) and integrated member of the \emph{Laboratory of Robotics and
Systems Engineering} (LARSyS).

Alexandre Bernardino started his pedagogical activity in 1995 and since
then has taught \textbf{21 different disciplinary curricular units} of
the BSc, MSc, and PhD degrees at IST/UL, on the broad areas of Robotics,
Control, Signal Processing and Artificial Intelligence. He has
supervised and co-supervised \textbf{9 postdoctoral fellows} and
graduated \textbf{22 PhD students} and more than \textbf{210 MSc
students}. He has hosted \textbf{23 visiting students} in their
internships at ISR and supervised \textbf{84 research fellowships}
integrated in research projects.

The research interests of Alexandre Bernardino span the areas of
Bioinspired and Learning Robots, Vision for Decision and Control, Object
Grasping and Manipulation, and Human-Robot Interaction. He has published
numerous peer-review articles in top international publications:
\textbf{112 in journals}, \textbf{207 in conferences}, and \textbf{11 in
book chapters}. He has \textbf{463 co-authors} in scientific
publications.

He participated in more than \textbf{40 national and international
research projects}, being the global coordinator in \textbf{8 of them}.
He was the Technical Manager of the EU Project Handle, and the local
coordinator of the IST team participation on the ERC Project ORIENT led
by Prof. John Van Opstal of the Radboud University, Nijmegen.

Beyond the significant pedagogical and scientific achievements,
Alexandre Bernardino has contributed to numerous managerial,
professional, knowledge transfer and outreach initiatives. He is a
current member of the \textbf{executive board of ISR},
\textbf{coordinator of the LARSyS Thematic Line} INTERACTION,
\textbf{adjunct coordinator of the scientific area} of Systems, Decision
and Control, and \textbf{adjunct coordinator of the Computer and Robot
Vision Lab} (VisLab) of IST. He was member of the \textbf{executive
board of the ECE department} from 2010 to 2016, \textbf{chair of the
IEEE Portuguese RAS Chapter} from 2014 to 2019, and \textbf{joint
coordinator of the Gender Balance interest group at IST from 2020 to
2026}. He was the \textbf{founder of a spin-off company}, and author of
the disclosure of \textbf{3 inventions}. He has co-supervised several
students in \textbf{theses with industry} and has established several
\textbf{contracts and protocols with industrial partners}.
